\providecommand{\Cscen}[1]{\mathcal C_{#1}}
\providecommand{\Csq}{\mathcal C_{\square}}
\providecommand{\Ccyc}{\mathcal C_{C_m}}
\providecommand{\Real}{\operatorname{Re}}
\providecommand{\sgn}{\operatorname{sgn}}

\section{Cycles}\label{sec:cycles}

Section~\ref{sec:pairsource} writes the three finite tests for an arbitrary pair-source scenario. For every cycle and every order we construct a strictly positive law that the strongest of them accepts and that no model produces, then measure how far the accepted laws can sit from compatibility.

Throughout, the $m$-cycle $C_m$, $m\ge3$, has observed variables $O_0,\dots,O_{m-1}$ and sources $S_0,\dots,S_{m-1}$, indices read modulo $m$, with $O_v$ a child of $S_{v-1}$ and $S_v$; the copied observations of Definition~\ref{def:ginflation} are written $O_v^{ij}$, children of $S_{v-1,i}$ and $S_{v,j}$. The triangle is $C_3$, with $\Ctri=\Cscen{C_3}$, read as $A=O_0$, $B=O_1$, $C=O_2$ and $X=S_0$, $Y=S_1$, $Z=S_2$; the square is $C_4$, written $\square$, read as $A(X,W)$, $B(X,Y)$, $C(Y,Z)$, $D(Z,W)$ with $X=S_0$, $Y=S_1$, $Z=S_2$, $W=S_3$.

\subsection{Parity witnesses on every cycle}\label{sec:cycleparity}

\begin{lemma}[Auxiliary parity density]\label{lem:parityaux}
For $N\ge1$ and $q\ge0$ let
\begin{equation}\label{eq:Hdensity}
H_{N,q}(s)=2^{-N}\sum_{\substack{S\subseteq[N]\\|S|\text{ even}}}(-q)^{|S|/2}\prod_{w\in S}s_w,
\qquad s\in\{-1,1\}^N .
\end{equation}
Then $H_{N,q}$ is normalized, its characters are
\begin{equation}\label{eq:Hmoments}
\E_{H_{N,q}}\Bigl[\prod_{w\in S}s_w\Bigr]=
\begin{cases}
0,&|S|\text{ odd},\\
(-q)^{|S|/2},&|S|\text{ even},
\end{cases}
\end{equation}
and if $N^2q\le1/4$ then $H_{N,q}(s)\ge\tfrac23\,2^{-N}>0$ for every $s$.
\end{lemma}

\begin{proof}
Normalization and \eqref{eq:Hmoments} are orthogonality of the characters of $\{-1,1\}^N$ under the uniform measure. For positivity, $\binom N{2j}\le N^{2j}$, so
\[
\bigl|2^NH_{N,q}(s)-1\bigr|\le\sum_{j\ge1}\binom N{2j}q^j\le\sum_{j\ge1}(N^2q)^j\le\sum_{j\ge1}4^{-j}=\tfrac13 . \qedhere
\]
\end{proof}

For $F\subseteq\{0,\dots,m-1\}$ let $\partial F=\{w:|F\cap\{w,w+1\}|=1\}$ be the set of edges of $C_m$ incident to an odd number of vertices of $F$. Every vertex meets two edges, so $|\partial F|$ is even.

\begin{definition}[Cycle parity family]\label{def:cycletarget}
For $m\ge3$ and $q\ge0$ let $P_{m,q}$ be the signed measure on $\{-1,1\}^m$ with density $2^{-m}\sum_F(-q)^{|\partial F|/2}\prod_{v\in F}o_v$ against the counting measure, so that it has total mass $1$ and
\begin{equation}\label{eq:cycletarget}
\E_{P_{m,q}}\Bigl[\prod_{v\in F}O_v\Bigr]=(-q)^{|\partial F|/2}\qquad\text{for every }F .
\end{equation}
\end{definition}

If $m^2q\le1/4$ then $P_{m,q}$ is a probability law, being the law of $(\sigma_{v-1}\sigma_v)_{v=0}^{m-1}$ for $\sigma\sim H_{m,q}$. The exact ranges in the two cases we use are wider.

Taking $F$ the whole vertex set gives $\E[\prod_vO_v]=1$, so wherever $P_{m,q}$ is a probability law the product of all $m$ signs is $1$ almost surely; taking $F$ a contiguous arc gives mean $-q$. For $m=3$ the family is the parity-even law with observed means $x,y,z$,
\begin{equation}\label{eq:paritylaw}
\Pi(x,y,z)(\alpha,\beta,\gamma)=\tfrac14\bigl(1+x\alpha+y\beta+z\gamma\bigr)\ \text{ on }\alpha\beta\gamma=1,\qquad P_{3,q}=\Pi(-q,-q,-q),
\end{equation}
which is a probability law exactly for $0\le q\le1/3$; and for $m=4$ it is the square parity family $P_q:=P_{4,q}$ with atoms
\begin{equation}\label{eq:squaretarget}
\tfrac18(1-6q+q^2)\ \text{ at }0000,\qquad
\tfrac18(1-q^2)\ \text{ at }0011,0110,1001,1100,\qquad
\tfrac18(1+q)^2\ \text{ at }0101,1010,1111,
\end{equation}
zero on odd parity, in the convention $\text{false}=0\leftrightarrow+1$. It is a law exactly for $0\le q\le3-2\sqrt2$.

Fix $t\ge1$, put $N=mt$ and draw auxiliary signs $s=(s_{v,i})_{v\in\{0,\dots,m-1\},\,i\in[t]}$ from $H_{N,q}$. Define every copied observation by
\begin{equation}\label{eq:cyclewitness}
O_v^{ij}=s_{v-1,i}\,s_{v,j}.
\end{equation}
For a set $F$ of copied observations let $\partial F$ denote the set of copied sources occurring an odd number of times in $\prod_{O\in F}O$; each copied observation contributes two incidences, so $|\partial F|$ is even, and \eqref{eq:Hmoments} gives
\begin{equation}\label{eq:masterchar}
\E\Bigl[\prod_{O\in F}O\Bigr]=(-q)^{|\partial F|/2}.
\end{equation}

\begin{lemma}[Boundaries of prescribed characters]\label{lem:boundary}
Let $F_1,\dots,F_r$ be pairwise ancestrally independent injectable sets of copied observations of $C_m$ and let $U_i\subseteq F_i$. Then each $\partial U_i$ has even cardinality and contains at most one copy of each source, so it is empty or meets at least two source families; the sets $\partial U_1,\dots,\partial U_r$ are pairwise disjoint; and $\partial(U_1\cup\dots\cup U_r)=\bigsqcup_i\partial U_i$, which again is empty or meets at least two source families. Moreover $\partial U_i$ is the image, under erasure of copy indices, of the boundary of the corresponding original observed set.
\end{lemma}

\begin{proof}
An injectable set contains at most one copied observation over each original vertex and its ancestry contains at most one copy of each original source, so $\partial U_i$ contains at most one coordinate from each family; being of even cardinality it is empty or has at least two elements in distinct families. Erasing copy indices carries the incidence count of $U_i$ to that of its image, which gives the last claim. Ancestrally independent sets have disjoint ancestries, hence disjoint coordinate sets and disjoint boundaries, and the boundary of a union of sets with disjoint ancestries is the disjoint union of the boundaries. If that union is nonempty then some $\partial U_i$ is nonempty and already meets two families.
\end{proof}

\begin{lemma}[The parity witness passes every prescription]\label{lem:cyclewitness}
Let $m\ge3$, $t\ge1$ and $0<q\le1/(4m^2t^2)$. The law of \eqref{eq:cyclewitness} is a strictly positive witness placing $P_{m,q}$ in $\IAI_t(C_m)=\Iexp_t(C_m)$.
\end{lemma}

\begin{proof}
Positivity of $H_{N,q}$ holds since $N^2q=m^2t^2q\le1/4$. The density \eqref{eq:Hdensity} is symmetric in the $N$ coordinates, so the witness is invariant under independent permutations of the copy indices of each source. By Lemma~\ref{lem:boundary} an injectable set $F$ has $\partial F$ equal to the boundary of its image, so \eqref{eq:masterchar} and \eqref{eq:cycletarget} give it the prescribed marginal of $P_{m,q}$. For pairwise ancestrally independent injectable $F_1,\dots,F_r$ and any $U_i\subseteq F_i$, Lemma~\ref{lem:boundary} and \eqref{eq:masterchar} give
\[
\E\Bigl[\prod_i\prod_{O\in U_i}O\Bigr]=(-q)^{\frac12\sum_i|\partial U_i|}=\prod_i\E\Bigl[\prod_{O\in U_i}O\Bigr],
\]
and characters determine a law on a finite sign cube, so the blocks are independent with the prescribed marginals. The diagonal rows are pairwise ancestrally independent injectable copied cycles, so the diagonal law is $P_{m,q}^{\otimes t}$. Lemma~\ref{lem:rootsink} supplies the recursive prescriptions.
\end{proof}

The witness is a law on the copied observations, not the output of an inflated model: under $H_{N,q}$ the auxiliary signs are correlated across families.

\begin{lemma}[Exact parity rigidity]\label{lem:parityrigidity}
Let $R\in\Ctri$ with $ABC=1$ almost surely. Then there are $\sigma,\tau,\upsilon\in[-1,1]$ with $(\E A,\E B,\E C)=(\sigma\tau,\sigma\upsilon,\tau\upsilon)$. In particular $\E A\cdot\E B\cdot\E C\ge0$.
\end{lemma}

\begin{proof}
Absorb the response seeds into incident sources, so that $A=A(X,Z)$, $B=B(X,Y)$, $C=C(Z,Y)$ are deterministic signs. For $\mu_Y$-almost every $y$ the identity $A(X,Z)B(X,y)C(Z,y)=1$ holds for $\mu_X\otimes\mu_Z$-almost every $(X,Z)$; fix such a $y_0$ and put $S(X)=B(X,y_0)$, $T(Z)=C(Z,y_0)$, so that $A=ST$ almost surely. Substituting in $ABC=1$ gives $\bigl(S(X)B(X,Y)\bigr)\bigl(T(Z)C(Z,Y)\bigr)=1$ almost surely. Write $B'=SB$ and $C'=TC$. Conditionally on $Y$ these are independent signs, and they agree almost surely, so each equals a function $U(Y)$ almost surely. Hence $A=ST$, $B=SU$, $C=TU$ almost surely, and independence of the three sources gives the means with $\sigma=\E S$, $\tau=\E T$, $\upsilon=\E U$. Their product is $(\sigma\tau\upsilon)^2$.
\end{proof}

Rigidity of the parity-perfect (parity token-counting) triangle correlations is prior work: the exact statement is due to Boreiri, Girardin, Ulu, Lipka-Bartosik, Brunner and Sekatski \cite{BoreiriEtAl2023}, and Boreiri, Ulu, Brunner and Sekatski \cite{BoreiriEtAl2025} give a noise-robust form (a compatible law with parity error $\epsilon$ is within $3\epsilon$ of a parity-perfect compatible law); Lemma~\ref{lem:quantrigidity} is a moment form of that robustness. Renou and Beigi \cite{RenouBeigi2022PRA,RenouBeigi2022PRL} prove rigidity of token-counting distributions on every ring scenario with bipartite sources; their statements concern integer token counts rather than parity, so they do not contain the lemma below, but they are the prior rigidity results on cycles. We give the two proofs because we need the statements for arbitrary latent alphabets and stochastic responses, and because the constant of Lemma~\ref{lem:quantrigidity} enters the distance bounds below.

\begin{lemma}[Quantitative parity rigidity]\label{lem:quantrigidity}
Let $R\in\Ctri$ and $W=\E_R[ABC]$. Then
\[
\max\{\E_RA,\ \E_RB,\ \E_RC\}\ \ge\ -4(1-W).
\]
\end{lemma}

\begin{proof}
Let $a(x,z)$, $b(x,y)$, $c(z,y)$ be the conditional means of $A$, $B$, $C$ given the sources; they take values in $[-1,1]$, and since the three responses use independent private randomness, $W=\E[abc]$ and $\E_RA=\E[a]$, with the expectations over the independent sources $X,Y,Z$. We use two inequalities: for $p,r\in[-1,1]$,
\begin{equation}\label{eq:pr}
|p-r|\le1-pr,
\end{equation}
since for $p\ge r$ the difference $(1-pr)-(p-r)=(1-p)(1+r)$ is nonnegative, and for $u,w\in[-1,1]$,
\begin{equation}\label{eq:uw}
|u+w|\ge2uw,
\end{equation}
since for $uw\le0$ the right side is nonpositive and for $u,w$ of the same sign $|u|+|w|-2|u||w|=|u|(1-|w|)+|w|(1-|u|)\ge0$.

Choose $y_0$ with $\E_{X,Z}[1-a(X,Z)b(X,y_0)c(Z,y_0)]\le1-W$; it exists because the average over $Y$ of the left side is $1-W$. Put $S(X)=b(X,y_0)$ and $T(Z)=c(Z,y_0)$. By \eqref{eq:pr} with $p=a$, $r=ST$,
\[
\E|a-ST|\le\E[1-aST]\le1-W,\qquad\text{so}\qquad |\E[a]-\E[S]\,\E[T]|\le1-W,
\]
using the independence of $X$ and $Z$. Next let $u(y)=\E_X[S(X)b(X,y)]$ and $w(y)=\E_Z[T(Z)c(Z,y)]$, both in $[-1,1]$, and $\upsilon(y)=\sgn(u(y)+w(y))$. By \eqref{eq:uw},
\begin{multline*}
\E_Y[1-\upsilon u]+\E_Y[1-\upsilon w]=\E_Y[2-|u+w|]\le2-2\E_Y[uw]\\
=2-2\E[STbc]\le2-2\bigl(W-\E|a-ST|\bigr)\le4(1-W),
\end{multline*}
where $\E_Y[uw]=\E[S(X)T(Z)b(X,Y)c(Z,Y)]$ by independence of the sources and $\E[STbc]\ge\E[abc]-\E|ST-a|$ because $|bc|\le1$. Hence $\E_Y[1-\upsilon u]\le4(1-W)$ and $\E_Y[1-\upsilon w]\le4(1-W)$. By \eqref{eq:pr} with $p=b(X,Y)$ and $r=S(X)\upsilon(Y)$,
\[
\E|b-S\upsilon|\le\E[1-bS\upsilon]=\E_Y[1-\upsilon u]\le4(1-W),\qquad\text{so}\qquad|\E[b]-\E[S]\,\E[\upsilon]|\le4(1-W),
\]
and likewise $|\E[c]-\E[T]\,\E[\upsilon]|\le4(1-W)$. Write $\sigma=\E[S]$, $\tau=\E[T]$, $\nu=\E[\upsilon]$. If all three means were below $-4(1-W)$, then $\sigma\tau\le\E[a]+(1-W)<0$, $\sigma\nu<0$ and $\tau\nu<0$, whose product $(\sigma\tau\nu)^2$ is nonnegative, a contradiction.
\end{proof}
For $W=1$ the lemma says that some mean is nonnegative, which is Lemma~\ref{lem:parityrigidity} again. A compatible law at total-variation distance $\delta$ from a parity-perfect law has $1-W\le2\delta$, so the lemma also replaces the repair step of \cite{BoreiriEtAl2025}: we construct no repaired model, and the bound holds for arbitrary latent alphabets and stochastic responses.

\begin{lemma}[Distance of the cycle family]\label{lem:cycledistance}
Let $m\ge3$ and let $q>0$ be such that $P_{m,q}$ is a probability law. Then $d_{\TV}(P_{m,q},\Ccyc)\ge q/10$; in particular $P_{m,q}\notin\Ccyc$.
\end{lemma}

\begin{proof}
Partition $\{0,\dots,m-1\}$ into three nonempty contiguous arcs $\Lambda_1,\Lambda_2,\Lambda_3$ and let $\varphi(o)=(V_1,V_2,V_3)$ with $V_r=\prod_{v\in\Lambda_r}o_v$. If $R\in\Ccyc$, then $\varphi_*R\in\Ctri$: the three sources on the arc boundaries are independent, each arc product is a function of the two boundary sources adjacent to its arc together with the sources internal to that arc, and those internal sources can be absorbed as private randomness of the arc. Each $\partial\Lambda_r$ has two elements, and $\partial(\Lambda_1\cup\Lambda_2)=\partial\Lambda_3$, so \eqref{eq:cycletarget} gives
\[
\varphi_*P_{m,q}=\Pi(-q,-q,-q).
\]

Let $R\in\Ccyc$ and $\delta=d_{\TV}(P_{m,q},R)$. Total variation does not increase under $\varphi$, so $Q=\varphi_*R\in\Ctri$ satisfies $d_{\TV}(\Pi(-q,-q,-q),Q)\le\delta$. Since $|\E_Pf-\E_Qf|\le2d_{\TV}(P,Q)$ for $|f|\le1$, the parity $W=\E_Q[V_1V_2V_3]$ satisfies $1-W\le2\delta$ and each mean $\E_QV_r$ is at most $-q+2\delta$. Lemma~\ref{lem:quantrigidity} gives $-q+2\delta\ge-4(1-W)\ge-8\delta$, that is $\delta\ge q/10$.
\end{proof}

\begin{theorem}[Cycles at every order]\label{thm:cycle}
Let $m\ge3$ and $t\ge1$, put
\[
q=\frac1{4m^2t^2},\qquad \eta=\frac q{32m},
\]
and let $\widetilde P_{m,q}=T_\eta^{\otimes m}P_{m,q}$, where $T_\eta$ flips a sign with probability $\eta$. Then $\widetilde P_{m,q}$ is a strictly positive rational law with
\[
\widetilde P_{m,q}\in\Iexp_t(C_m)=\IAI_t(C_m),\qquad
d_{\TV}(\widetilde P_{m,q},\Ccyc)\ \ge\ \frac{11q}{160}=\frac{11}{640\,m^2t^2}>0 .
\]
Consequently no finite order of any of the three hierarchies characterizes $\Ccyc$.
\end{theorem}

\begin{proof}
Apply $T_\eta$ independently to every copied observation of the witness of Lemma~\ref{lem:cyclewitness}. The resulting law is strictly positive, still invariant under the copy-index permutations, and its characters are those of \eqref{eq:masterchar} multiplied by $(1-2\eta)^{|F|}$. The same factor multiplies the target characters, and $|F|$ adds over disjoint blocks, so every injectable marginal and every ancestrally independent product of Lemma~\ref{lem:cyclewitness} persists with $\widetilde P_{m,q}$ in place of $P_{m,q}$, as does the diagonal law. Lemma~\ref{lem:rootsink} again supplies the recursive prescriptions.

For the distance, $d_{\TV}(\widetilde P_{m,q},P_{m,q})\le m\eta=q/32$ by a union bound over the flipped coordinates, and Lemma~\ref{lem:cycledistance} gives $d_{\TV}(\widetilde P_{m,q},\Ccyc)\ge q/10-q/32=11q/160$. Since $t$ was arbitrary and $\widetilde P_{m,q}\notin\Ccyc$, the feasible set at order $t$ is strictly larger than $\Ccyc$ for every $t$.
\end{proof}

For $m=3$ this gives a worst-case distance of order $t^{-2}$ for the triangle, in place of the order $t^{-3}$ that the family of Section~\ref{sec:main} supplies for Question~\ref{q:hausdorff}.

\subsection{A corrected density and a linear lower bound}\label{sec:corrected}

The density \eqref{eq:Hdensity} fixes every auxiliary character, including characters that no prescription refers to. Changing those characters restores positivity at a larger $q$.

\begin{lemma}[Phase telescoping]\label{lem:telescope}
For real $\theta_1,\dots,\theta_m$,
\[
1-\cos\Bigl(\sum_{g=1}^m\theta_g\Bigr)\ \le\ m\sum_{g=1}^m\bigl(1-\cos\theta_g\bigr).
\]
\end{lemma}

\begin{proof}
$1-\cos\phi=\tfrac12|1-e^{i\phi}|^2$, and
$1-\prod_ge^{i\theta_g}=\sum_g\bigl(\prod_{g'<g}e^{i\theta_{g'}}\bigr)(1-e^{i\theta_g})$,
so $|1-\prod_ge^{i\theta_g}|\le\sum_g|1-e^{i\theta_g}|$. Cauchy--Schwarz gives $\bigl(\sum_g a_g\bigr)^2\le m\sum_ga_g^2$.
\end{proof}

\begin{lemma}[Corrected density]\label{lem:corrected}
Fix $m\in\{3,4\}$, $t\ge1$ and $0<q\le1/(16t)$, and set $c=4$ for $m=3$ and $c=5$ for $m=4$. On $m$ families of $t$ signs $s_{g,i}$ put
\begin{equation}\label{eq:Wdensity}
f_g=\prod_{i=1}^t\bigl(1+i\sqrt q\,s_{g,i}\bigr),\qquad R=(1+q)^{t/2},\qquad
W(s)=\Real\prod_{g=1}^mf_g+c\sum_{g=1}^m\bigl(1-\Real f_g\bigr).
\end{equation}
Then $W$ is a real polynomial in $q$ with rational coefficients, $W\ge3/5$ for $m=3$ and $W\ge1/3$ for $m=4$, the law $\rho=2^{-mt}W$ is a probability law, and its characters are
\begin{equation}\label{eq:Wmoments}
\E_\rho\Bigl[\prod_{w\in S}s_w\Bigr]=
\begin{cases}
0,&|S|\text{ odd},\\
(1-c)(-q)^{|S|/2},&S\ne\varnothing\text{ even, inside one family},\\
(-q)^{|S|/2},&|S|\text{ even, meeting at least two families}.
\end{cases}
\end{equation}
\end{lemma}

\begin{proof}
Taking real parts keeps only even powers of $i\sqrt q$, so $W$ is a real polynomial in $q$. Since $|1+i\sqrt q\,s|=(1+q)^{1/2}$ we may write $f_g=Re^{i\theta_g}$, and then
\[
\begin{aligned}
W&=R^m\cos\Bigl(\sum_g\theta_g\Bigr)+c\sum_g\bigl(1-R\cos\theta_g\bigr)\\
&=R^m-cm(R-1)-R^m\Bigl(1-\cos\sum_g\theta_g\Bigr)+cR\sum_g\bigl(1-\cos\theta_g\bigr).
\end{aligned}
\]
Lemma~\ref{lem:telescope} gives
\[
W\ \ge\ R^m-cm(R-1)+\bigl(cR-mR^m\bigr)\sum_g\bigl(1-\cos\theta_g\bigr).
\]
From $tq\le1/16$ we get $R^2=(1+q)^t\le\sum_{j\ge0}(tq)^j\le16/15$ and $R-1=(R^2-1)/(R+1)\le1/30$. For $m=3$, $c=4$: $3R^3\le4R$ because $3R^2\le16/5\le4$, and $W\ge1-12/30=3/5$. For $m=4$, $c=5$: $4R^4\le5R$ because $4R^3\le4(16/15)^{3/2}<5$, and $W\ge1-20/30=1/3$.

For the characters, work under the uniform law on the sign cube, where the families are independent and $\E f_g=1$. For $S$ with parts $S_g$ in the families,
\[
\E\Bigl[\prod_gf_g\prod_{w\in S}s_w\Bigr]=\prod_g\E\Bigl[f_g\prod_{w\in S_g}s_w\Bigr]=(i\sqrt q)^{|S|},
\]
whose real part is $0$ for odd $|S|$ and $(-q)^{|S|/2}$ for even $|S|$. Also $\E[(1-\Real f_g)\prod_{w\in S}s_w]$ vanishes unless $S$ is a nonempty subset of family $g$, in which case it equals $-\Real(i\sqrt q)^{|S|}$. Adding the contributions gives \eqref{eq:Wmoments}; the case $S=\varnothing$ gives $\E_\rho1=1$.
\end{proof}

\begin{theorem}[The square at $q=\Theta(1/t)$]\label{thm:squarelinear}
For every $t\ge1$ and every $0<q\le1/(16t)$, the square parity law $P_q$ of \eqref{eq:squaretarget} lies in $\IAI_t(\square)=\Iexp_t(\square)$.
\end{theorem}

\begin{proof}
Take $m=4$ and $c=5$ in Lemma~\ref{lem:corrected}, name the four families $x,y,z,w$ after the sources $X,Y,Z,W$, and define every copied observation by
\[
A^{il}=x_iw_l,\qquad B^{ij}=x_iy_j,\qquad C^{jk}=y_jz_k,\qquad D^{kl}=z_kw_l,
\]
that is by \eqref{eq:cyclewitness} for $m=4$. The density \eqref{eq:Wdensity} is symmetric within each family, so the law is invariant under independent permutations of the copy indices of each source. By Lemma~\ref{lem:boundary} every boundary occurring in an AI prescription, and every boundary of a subset of one of its blocks, is empty or meets at least two families; by \eqref{eq:Wmoments} such characters take the value $(-q)^{|\partial\cdot|/2}$, unchanged from \eqref{eq:masterchar}. The computation of Lemma~\ref{lem:cyclewitness} therefore applies verbatim: every injectable set has the marginal of $P_{4,q}=P_q$, ancestrally independent injectable blocks are independent, and the diagonal law is $P_q^{\otimes t}$. The moments that Lemma~\ref{lem:corrected} changed are the nonempty even characters inside a single family, and no prescription refers to one. Lemma~\ref{lem:rootsink} gives the recursive prescriptions.
\end{proof}

\begin{theorem}[The triangle at $q=\Theta(1/t)$]\label{thm:trianglelinear}
For every $t\ge1$ and every $0<q\le1/(16t)$, the parity-even law $\Pi(-q,-q,-q)$ of \eqref{eq:paritylaw} lies in $\IAI_t(\triangle)=\Iexp_t(\triangle)$.
\end{theorem}

\begin{proof}
Take $m=3$ and $c=4$ in Lemma~\ref{lem:corrected}, name the three families $x,y,z$ after the sources $X,Y,Z$, and define
\[
A^{ij}=x_iz_j,\qquad B^{ik}=x_iy_k,\qquad C^{jk}=z_jy_k,
\]
which is \eqref{eq:cyclewitness} for $m=3$ in the index convention of Section~\ref{sec:setup}. Positivity and normalization are Lemma~\ref{lem:corrected} with $W\ge3/5$. Symmetry within each family is again immediate. By Lemma~\ref{lem:boundary} every character prescribed by an injectable marginal, by an ancestrally independent product, or by the diagonal law has boundary empty or meeting at least two families, so \eqref{eq:Wmoments} returns $(-q)^{|\partial\cdot|/2}$ and the argument of Lemma~\ref{lem:cyclewitness} applies unchanged. The original copied triangle has $\E A=\E[x_iz_j]=-q$ and likewise $\E B=\E C=-q$, while $ABC=x_iz_jx_iy_kz_jy_k=1$, so its law is $\Pi(-q,-q,-q)=P_{3,q}$. Lemma~\ref{lem:rootsink} gives the recursive prescriptions.
\end{proof}
The exact certificates described in Section~\ref{sec:repro} check the square and triangle witnesses of the two theorems, with and without the flips of Corollaries~\ref{cor:squaredistance} and~\ref{cor:triangledistance}, at orders one and two for the square and one to three for the triangle: positivity of every atom, symmetry, the full diagonal law and every ancestral-independence prescription.


\subsection{Distance and the worst-case profile}\label{sec:worstcase}

For $H\in\{\mathrm{NW},\mathrm{AI},\mathrm{exp}\}$ and a pair-source scenario $G$ set
\begin{equation}\label{eq:Hprofile}
H^H_t(G)=\sup_{P\in\mathcal I^H_t(G)}d_{\TV}(P,\Cscen{G}),
\end{equation}
the worst-case total-variation distance from compatibility of a law the order-$t$ test accepts, with $d_{\TV}$ half the $\ell^1$ distance as in Section~\ref{sec:setup}.

\begin{lemma}[Max-moment inequality for the square]\label{lem:maxmoment}
Every $P\in\Csq$ in sign notation satisfies
\[
\max\{\E A,\ \E B,\ \E[AB]\}+4\,P(ABCD=-1)\ \ge\ 0 .
\]
\end{lemma}

\begin{proof}
Absorb the response seeds into incident sources, so $A=A(X,W)$, $B=B(X,Y)$, $C=C(Y,Z)$, $D=D(Z,W)$ are deterministic signs and $X,Y,Z,W$ are independent. Fix $z$ and write $g_z(Y)=C(Y,z)$, $h_z(W)=D(z,W)$ and
\[
\delta_z=\Pr\bigl(A(X,W)B(X,Y)\ne g_z(Y)h_z(W)\bigr),
\]
so that $\E_Z\delta_Z=P(ABCD=-1)=:\delta$. Put
\[
u_z(x)=\E_W\bigl[A(x,W)h_z(W)\bigr],\qquad v_z(x)=\E_Y\bigl[B(x,Y)g_z(Y)\bigr].
\]
Since $Y$ and $W$ are independent, $\E[ABg_zh_z\mid X=x]=u_z(x)v_z(x)$, and as $ABg_zh_z$ is a sign, $\E_X[1-u_zv_z]=2\delta_z$.

Let $f_z(x)=\sgn(u_z(x)+v_z(x))$, with $\sgn0=1$, and compare $A$ and $B$ with $A'_z=f_z(X)h_z(W)$ and $B'_z=f_z(X)g_z(Y)$. Then
\[
e_z:=\Pr(A\ne A'_z)+\Pr(B\ne B'_z)=\E_X\Bigl[1-\tfrac12|u_z+v_z|\Bigr]\le\E_X\bigl[1-u_zv_z\bigr]=2\delta_z,
\]
using $|u+v|\ge2uv$ on $[-1,1]^2$. Independence of $X,Y,W$ gives
\[
(\E A'_z)(\E B'_z)(\E[A'_zB'_z])=(\E f_z)^2(\E g_z)^2(\E h_z)^2\ \ge\ 0,
\]
so at least one of the three comparison moments is nonnegative. Each differs from the corresponding moment of $P$ by at most $2e_z$, since $|\E U-\E V|\le2\Pr(U\ne V)$ for signs and $\Pr(AB\ne A'_zB'_z)\le\Pr(A\ne A'_z)+\Pr(B\ne B'_z)$. Hence
\[
\max\{\E A,\E B,\E[AB]\}\ \ge\ -2e_z\ \ge\ -4\delta_z .
\]
The left side does not depend on $z$; averaging over $z$ finishes the proof.
\end{proof}

\begin{corollary}[Full-support survivors on the square]\label{cor:squaredistance}
$d_{\TV}(P_q,\Csq)\ge q/6$ for $0<q\le3-2\sqrt2$. With flip probability $\eta=q/64$ applied to each of the four bits and to every copied observation, $\widetilde P_q=T_\eta^{\otimes4}P_q$ is strictly positive, lies in $\IAI_t(\square)=\Iexp_t(\square)$ for $q\le1/(16t)$, and satisfies $d_{\TV}(\widetilde P_q,\Csq)\ge5q/48$. At $q_t=1/(16t)$,
\[
H^{\mathrm{exp}}_t(\square)=H^{\mathrm{AI}}_t(\square)\ \ge\ \frac5{768\,t}.
\]
\end{corollary}

\begin{proof}
If $R\in\Csq$ and $d=d_{\TV}(P_q,R)$, each of $\E_RA$, $\E_RB$, $\E_R[AB]$ is at most $-q+2d$ and the parity error of $R$ is at most $d$, since $P_q$ has none. Lemma~\ref{lem:maxmoment} gives $0\le-q+2d+4d$, so $d\ge q/6$. The noise argument of Theorem~\ref{thm:cycle} applies to the witness of Theorem~\ref{thm:squarelinear}, and $d_{\TV}(\widetilde P_q,P_q)\le4\eta=q/16$, so $d_{\TV}(\widetilde P_q,\Csq)\ge q/6-q/16=5q/48$. Substituting $q_t=1/(16t)$ gives $5/(768t)$, and Lemma~\ref{lem:rootsink} identifies the two suprema.
\end{proof}

\begin{corollary}[Full-support survivors on the triangle]\label{cor:triangledistance}
$d_{\TV}(\Pi(-q,-q,-q),\Ctri)\ge q/10$ for $0<q\le1/3$. With flip probability $\eta=q/192$ on each of the three bits and on every copied observation, the noisy target is strictly positive, lies in $\IAI_t(\triangle)=\Iexp_t(\triangle)$ for $q\le1/(16t)$, and is at distance at least $27q/320$ from $\Ctri$. At $q_t=1/(16t)$,
\[
H^{\mathrm{exp}}_t(\triangle)=H^{\mathrm{AI}}_t(\triangle)\ \ge\ \frac{27}{5120\,t}.
\]
\end{corollary}

\begin{proof}
The distance bound is Lemma~\ref{lem:cycledistance} for $m=3$, whose proof uses only Lemma~\ref{lem:quantrigidity} on the triangle itself. The three flips cost at most $3\eta=q/64$ in total variation, and $q/10-q/64=27q/320$. Feasibility is Theorem~\ref{thm:trianglelinear} with the noise argument of Theorem~\ref{thm:cycle}. Substituting $q_t=1/(16t)$ gives $27/(5120t)$.
\end{proof}

\subsection{A convex-order upper bound}\label{sec:convexorder}

\begin{theorem}[Convex order along the diagonal]\label{thm:convexorder}
Let $G$ be a correlation scenario, that is, finitely many independent latent sources and observed variables with finite alphabets, each a stochastic function of the sources it reads, with $K$ joint outcomes; the order-$n$ inflation and the set $\INW_n(G)$ are defined as in Section~\ref{sec:pairsource} with sources in place of edges. Let $P\in\INW_n(G)$ with witness $\Gamma$. For a deterministic table $\omega$ of copied observations let $q_\omega$ be the observed law obtained by sampling one copy index uniformly and independently for each source and reading $\omega$ at those indices. Then $q_\omega\in\Cscen{G}$ for every $\omega$, and for every convex $\Phi$ on the observed simplex
\begin{equation}\label{eq:convexorder}
\E_\Gamma\,\Phi(q_\omega)\ \le\ \E\,\Phi(\widehat P_n),
\end{equation}
where $\widehat P_n$ is the empirical law of $n$ independent draws from $P$. In particular
\begin{equation}\label{eq:convexl2}
\E_\Gamma\|q_\omega-P\|_2^2\ \le\ \frac{1-\|P\|_2^2}{n},
\qquad
H^{\mathrm{NW}}_n(G)\ \le\ \min\Bigl\{1,\ \frac{\sqrt{K-1}}{2\sqrt n}\Bigr\}.
\end{equation}
\end{theorem}

\begin{proof}
Fix $\omega$. Sampling an index $I_e\in[n]$ uniformly and independently for each source $e$ and letting each observed variable $v$ output $\omega_v\bigl((I_e)_{e\ni v}\bigr)$, where $e\ni v$ ranges over the sources $v$ reads, is a model of $G$ with independent sources $I_e$ and deterministic responses, so $q_\omega\in\Cscen{G}$.

Draw one uniform permutation $\pi_e$ of $[n]$ for each source, independently of each other and of $\omega\sim\Gamma$, and set
\[
T_\ell=\Bigl(\omega_v\bigl((\pi_e(\ell))_{e\ni v}\bigr)\Bigr)_{v},\qquad \ell\in[n],
\qquad \widehat P=\frac1n\sum_{\ell=1}^n\delta_{T_\ell}.
\]
For each fixed choice of permutations, the symmetry of $\Gamma$ carries $(T_1,\dots,T_n)$ to the diagonal rows, so $(T_1,\dots,T_n)\sim P^{\otimes n}$; this remains true after averaging over the permutations, so $\widehat P$ has the law of $\widehat P_n$. Conditionally on $\omega$ and on a fixed $\ell$, the indices $(\pi_e(\ell))_e$ are independent and uniform on $[n]$, so $\E[\delta_{T_\ell}\mid\omega]=q_\omega$ and hence $\E[\widehat P\mid\omega]=q_\omega$. Conditional Jensen gives $\Phi(q_\omega)\le\E[\Phi(\widehat P)\mid\omega]$, and averaging over $\omega$ gives \eqref{eq:convexorder}.

Take $\Phi(\mu)=\|\mu-P\|_2^2$. Then $\E\|\widehat P_n-P\|_2^2=\frac1n\sum_xP(x)(1-P(x))=(1-\|P\|_2^2)/n$, which is the first part of \eqref{eq:convexl2}. Some $\omega$ attains at most the mean, and its $q_\omega$ is compatible, so
\[
d_{\TV}(P,\Cscen{G})\le\tfrac12\|P-q_\omega\|_1\le\frac{\sqrt K}2\|P-q_\omega\|_2\le\frac{\sqrt K}2\sqrt{\frac{1-\|P\|_2^2}n}\le\frac{\sqrt{K-1}}{2\sqrt n},
\]
using $\|P\|_2^2\ge1/K$. Total variation never exceeds $1$.
\end{proof}

The sampling mechanism is that of the finite de Finetti estimate of \cite[Section 4.2 and eq.~(A7)]{NW2020}, which bounds the same quantity by $L(1-\|P\|_2^2)/n$ for $L$ sources through the collision probability of two index samples. Using the full diagonal law through a conditional expectation, rather than its first two degrees, removes the factor $L$. Section~\ref{sec:general} draws the rejecting-order consequence.

\begin{corollary}[Brackets for the square and the triangle]\label{cor:brackets}
For every $t\ge1$,
\[
\frac5{768\,t}\ \le\ H^{\mathrm{exp}}_t(\square)=H^{\mathrm{AI}}_t(\square)\ \le\ H^{\mathrm{NW}}_t(\square)\ \le\ \min\Bigl\{1,\frac{\sqrt{15}}{2\sqrt t}\Bigr\},
\]
\[
\frac{27}{5120\,t}\ \le\ H^{\mathrm{exp}}_t(\triangle)=H^{\mathrm{AI}}_t(\triangle)\ \le\ H^{\mathrm{NW}}_t(\triangle)\ \le\ \frac{\sqrt7}{2\sqrt t}.
\]
\end{corollary}

\begin{proof}
Combine Corollaries~\ref{cor:squaredistance} and~\ref{cor:triangledistance} with Theorem~\ref{thm:convexorder} at $K=16$ and $K=8$, and with Proposition~\ref{prop:gsound}.
\end{proof}

The two exponents do not meet. Whether $H^{\mathrm{NW}}_t(\square)=O(1/t)$ is open, and so is the corresponding question for the triangle; this is the form Question~\ref{q:hausdorff} takes once the lower bounds above replace the family bound.

\subsection{Order conversion on the square}\label{sec:conversion}

\begin{theorem}[NW and AI differ by a bounded change of order]\label{thm:orderconversion}
For every $t\ge2$,
\[
\INW_{\lfloor3t/2\rfloor}(\square)\ \subseteq\ \IAI_t(\square)=\Iexp_t(\square)\ \subseteq\ \INW_t(\square).
\]
\end{theorem}

\begin{proof}
The right inclusion is Proposition~\ref{prop:gsound} and the equality is Lemma~\ref{lem:rootsink}. For the left inclusion let $\Gamma$ be an order-$n$ Navascu\'es--Wolfe witness for $P$ with $n=\lfloor3t/2\rfloor$ and let $\Gamma'$ be its restriction to the copy indices $1,\dots,t$ of each source. Restriction preserves symmetry and the diagonal law of degree $t$, so $\Gamma'\in\INW_t$; we check the AI prescriptions.

Fix an AI set of $\Gamma'$ and decompose it into the connected components of its shared-parent graph, which are injectable and pairwise ancestrally independent. A component's ancestry uses two, three or four source types; the two-type components are single copied observations and correspond to edges of the four-cycle $X-Y-Z-W-X$ of source types. Any two subsets of a four-element set of sizes at least three, or of sizes three and two, intersect, so a component using at least three source types conflicts with every other component and needs a diagonal row of its own. Let $b$ be the number of these.

The remaining components form a multigraph on the source-type four-cycle, with multiplicities $n_A,n_B,n_C,n_D$ on the edges $A,B,C,D$. Opposite edges are disjoint, so the colours $\max(n_A,n_C)$ for the pair $\{A,C\}$ and $\max(n_B,n_D)$ for the pair $\{B,D\}$ suffice, and their sum equals the maximum source-type degree $\Delta$: if $n_A\ge n_C$ and $n_B\ge n_D$ the sum is $n_A+n_B=\deg X$, and $\deg X$ dominates the other three degrees. So $b+\Delta$ rows suffice.

Choose a source type of residual degree $\Delta$, say $X$, together with its two neighbours $Y$ and $W$ on the cycle. A component using at least three source types omits at most one, so it uses at least two of $X,Y,W$; each of the $\Delta$ two-type components incident to $X$ uses exactly two of them. The components have pairwise disjoint copied ancestries, and the three families $X,Y,W$ supply $3t$ copies in $\Gamma'$, so $2b+2\Delta\le3t$ and therefore $b+\Delta\le\lfloor3t/2\rfloor=n$.

Assign each component to the diagonal row of its colour. Components sharing a row use disjoint source types, and distinct components use distinct copies within each family, so independent permutations of the copy indices of the four sources place all components on their assigned rows simultaneously. Under $\Gamma$ the rows are independent with law $P$. Within one row the components have disjoint original source types, so Lemma~\ref{lem:sourcedisjoint}, available because $n\ge2$, makes $P$ factor across them. Hence the joint law of the placed components is the product of their $P$-marginals. Undoing the permutations, the AI set of $\Gamma'$ has exactly the prescribed product.
\end{proof}

The factor $3/2$ is the cost of placing all components on diagonal rows at once; ancestrally disjoint collections of $\lfloor3t/2\rfloor$ injectable components whose source-type conflict graph is complete show that this method cannot do better. Hence $H^{\mathrm{NW}}_{\lfloor3t/2\rfloor}(\square)\le H^{\mathrm{AI}}_t(\square)\le H^{\mathrm{NW}}_t(\square)$, so the two hierarchies have the same polynomial decay exponent if either has one.

\subsection{Low-order thresholds}\label{sec:lowthresholds}

On the parity family we settle the first ten orders exactly. The reduction below turns feasibility for a parity-perfect target into a finite linear program in the numbers of negative auxiliary signs; we state it for the square, where it also covers the Navascu\'es--Wolfe test.

\begin{lemma}[Count-moment reduction]\label{lem:countmoment}
Let $0\le q\le3-2\sqrt2$ and $t\ge1$. Then $P_q\in\IAI_t(\square)$ if and only if there are $w_a\ge0$, $a=(a_X,a_Y,a_Z,a_W)\in\{0,\dots,t\}^4$, with $\sum_aw_a=1$ and
\begin{equation}\label{eq:countmoment}
\sum_aw_a\prod_{g}K^{(t)}_{k_g}(a_g)=(-q)^{S/2},\qquad S=\sum_gk_g,
\end{equation}
for exactly those degree tuples $k\in\{0,\dots,t\}^4$ with $S$ even and $2\max_gk_g\le S$, where
\[
K^{(t)}_k(a)=\binom tk^{-1}\sum_j(-1)^j\binom aj\binom{t-a}{k-j}
\]
is the normalized Krawtchouk polynomial.
\end{lemma}

\begin{proof}[Proof sketch]
Every copied square is injectable, so under any witness it has law $P_q$ and hence parity $+1$ almost surely. Parity forcing across the copy indices then gives sign potentials: from $A^{il}B^{ij}C^{jk}D^{kl}=1$ for all indices, the product $B^{ij}C^{jk}$ depends only on $(i,k)$, so writing $M^{ik}$ for it and $x_i=M^{i1}$, $y_j=C^{j1}$ gives $B^{ij}=x_iy_j$ and, by the same argument on the other pairs, every supported array is of the form $A^{il}=x_iw_l$, $B^{ij}=x_iy_j$, $C^{jk}=y_jz_k$, $D^{kl}=z_kw_l$. The potential is unique up to the simultaneous flip of all four families, so averaging over the two gauges and over independent permutations of the copy indices reduces a witness to a distribution on the four counts $a_g$ of negative signs. Under that symmetrization the average of a character with $k_g$ distinct indices in family $g$ is $\prod_gK^{(t)}_{k_g}(a_g)$, by the definition of the Krawtchouk polynomial as the mean of a $k$-element sign product drawn without replacement.

The degrees that the AI prescriptions constrain are exactly those in the statement. Necessity: by Lemma~\ref{lem:boundary} a nonempty prescribed boundary meets at least two source families, so no single family can carry more than half of the total degree, and the total degree is even. Sufficiency: a balanced degree tuple can be realized by ancestrally disjoint paths in the source four-cycle, matching opposite-type terminals through intermediate copies and then adjacent terminals directly, so the corresponding character is prescribed. The prescribed value is $(-q)^{S/2}$ by \eqref{eq:cycletarget}.
\end{proof}

\begin{lemma}[The Navascu\'es--Wolfe form]\label{lem:countmomentNW}
Under the hypotheses of Lemma~\ref{lem:countmoment}, $P_q\in\INW_t(\square)$ if and only if \eqref{eq:countmoment} holds for exactly the degree tuples in
\[
D_{\mathrm{NW}}(t)=\{b_1+\dots+b_t:\ b_r\in\{0,1\}^4\text{ of even weight}\}\subseteq D_{\mathrm{AI}}(t),
\]
where $D_{\mathrm{AI}}(t)$ is the degree set of Lemma~\ref{lem:countmoment}. The same reductions hold for the triangle family $\Pi(-q,-q,-q)$ with three families of signs and the potentials $A^{ij}=x_iz_j$, $B^{ik}=x_iy_k$, $C^{jk}=z_jy_k$.
\end{lemma}

\begin{proof}[Proof sketch]
Every copied square is carried to the first diagonal row by independent copy-index permutations, so under a symmetric witness with diagonal law $P_q^{\otimes t}$ it has law $P_q$ and parity $+1$; the potential representation and the symmetrization then proceed as in Lemma~\ref{lem:countmoment}. The diagonal law prescribes exactly the characters $\prod_r\prod_{v\in F_r}O_v^{(r)}$ with one vertex set $F_r$ per diagonal row; in potential form row $r$ contributes the indicator of $\partial F_r$ to the degree tuple, and the boundaries of vertex sets of a connected graph are exactly its even-weight edge sets. Conversely a distribution on counts with the prescribed moments yields a symmetric witness with the diagonal law by the construction in the proof of Lemma~\ref{lem:countmoment}. For the triangle, $A^{ij}B^{ik}C^{jk}=1$ at $j=1$ gives $B^{ik}=A^{i1}C^{1k}$, so $x_i=A^{i1}$ and $y_k=C^{1k}$ serve as potentials for $B$, and $z_j=A^{1j}x_1$ completes the representation, unique up to a global flip.
\end{proof}

\begin{lemma}[Monotonicity along the parity direction]\label{lem:paritymonotone}
Let $H\in\{\mathrm{NW},\mathrm{AI},\mathrm{exp}\}$, $G\in\{\triangle,\square\}$, and $0\le q'\le q$ in the range where the parity target is a law. If the parity target with parameter $q$ lies in $I^H_t(G)$, so does the one with parameter $q'$.
\end{lemma}

\begin{proof}
By Lemmas~\ref{lem:countmoment} and~\ref{lem:countmomentNW} we may take a witness to be a law on sign potentials. Multiply every potential sign by an independent $\pm1$ variable of mean $\rho=\sqrt{q'/q}$. A character whose boundary has $S$ elements is multiplied by $\rho^S$, so its value $(-q)^{S/2}$ becomes $(-q')^{S/2}$, while symmetry and the block structure of every prescription are unchanged. Lemma~\ref{lem:rootsink} covers the recursively expressible test.
\end{proof}

By Lemma~\ref{lem:paritymonotone} the set of parameters accepted at a given order is an interval $[0,q^H_t]$, and we bracket $q^H_t$ by a rational primal witness of Lemmas~\ref{lem:countmoment} and~\ref{lem:countmomentNW} at a point below it and rational dual certificates above it.

\begin{proposition}[The parity thresholds at orders up to ten]\label{prop:paritythresholds}
For the square and the triangle, the thresholds $q^{\mathrm{AI}}_t=q^{\mathrm{exp}}_t$ and $q^{\mathrm{NW}}_t$ satisfy the bounds of Table~\ref{tab:thresholds} for $1\le t\le10$, each lower bound being a rational witness and each upper bound a dual certificate valid on the whole interval up to the endpoint of the family. In particular:
\begin{enumerate}[label=(\roman*),nosep]
\item $q^{\mathrm{AI}}_t<q^{\mathrm{NW}}_t$ at every even $t\le10$, so $\IAI_t\subsetneq\INW_t$ there on both scenarios;
\item the two brackets coincide at every odd $t\le9$;
\item $q^{\mathrm{NW}}_{t+1}<q^{\mathrm{AI}}_t$ for every $t\le9$ except $t=3$, where $q^{\mathrm{NW}}_3$ and $q^{\mathrm{NW}}_4$ share the same bracket;
\item on the square, $q^{\mathrm{AI}}_2$ is the zero $r$ of $g(q)=q^3-33q^2+27q-3$ and $P_q\in\INW_2(\square)$ for every $q\le0.17156287$; on the triangle, $q^{\mathrm{AI}}_2=2-\sqrt3$, $q^{\mathrm{NW}}_3=q^{\mathrm{NW}}_4=q^{\mathrm{AI}}_3=1/5$ and $q^{\mathrm{AI}}_5=q^{\mathrm{NW}}_5=1/7$.
\end{enumerate}
Consequently $\INW_3(\square)\subsetneq\IAI_2(\square)=\Iexp_2(\square)\subsetneq\INW_2(\square)$.
\end{proposition}

\begin{table}
\centering
\small
\begin{tabular}{r@{\quad}ll@{\qquad}ll}
\toprule
& \multicolumn{2}{c}{square} & \multicolumn{2}{c}{triangle}\\
$t$ & $q^{\mathrm{AI}}_t$ & $q^{\mathrm{NW}}_t$ & $q^{\mathrm{AI}}_t$ & $q^{\mathrm{NW}}_t$\\
\midrule
1 & $3-2\sqrt2$ & $3-2\sqrt2$ & $1/3$ & $1/3$\\
2 & $0.1324743290$ & $\ge0.1715628752$ & $2-\sqrt3$ & $1/3$\\
3 & $0.0950228950$ & $0.0950228950$ & $1/5$ & $1/5$\\
4 & $0.0835279580$ & $0.0950228950$ & $0.1730047060$ & $1/5$\\
5 & $0.0657380020$ & $0.0657380020$ & $1/7$ & $1/7$\\
6 & $0.0603176600$ & $0.0645356170$ & $0.1283741760$ & $0.1381856900$\\
7 & $0.0502565110$ & $0.0502565110$ & $0.1092774840$ & $0.1092774840$\\
8 & $0.0470271490$ & $0.0494541280$ & $0.1011416960$ & $0.1070388380$\\
9 & $0.0406784540$ & $0.0406784540$ & $0.0888129680$ & $0.0888129680$\\
10 & $0.0384981990$ & $0.0400943600$ & $0.0833179940$ & $0.0870307060$\\
\bottomrule
\end{tabular}
\caption{Certified brackets for the parity thresholds. A decimal entry $x$ means $x\le q^H_t\le x+4\cdot10^{-9}$; the exact entries are the roots of the dual polynomials; the entry $\ge0.1715628752$ is a lower bound only, the family ending at $3-2\sqrt2\approx0.1715728753$.}\label{tab:thresholds}
\end{table}

\begin{proof}
All statements are exact certificates, described in Section~\ref{sec:repro}: for each order and hierarchy a nonnegative rational count witness at the lower endpoint, checked against every prescribed moment of Lemmas~\ref{lem:countmoment} and~\ref{lem:countmomentNW}, and a chain of rational dual functionals, each nonpositive on all count configurations, whose prescribed expectations are polynomials in $q$ certified positive by Sturm sequences on intervals covering everything above the upper endpoint. Lemma~\ref{lem:paritymonotone} turns the point witnesses into the intervals. The exact values arise where the dual polynomial factors: at square order two the dual is the functional $D$ with $\E D=(1+q)g(q)$, verified coefficient by coefficient; on the triangle the duals are $-(q+1)(q^2-4q+1)$ at order two and $(5q-1)(q+1)^2$, $(5q-1)(q+1)^3$ at orders three and four, and the order-five dual vanishes at $1/7$. The degree sets $D_{\mathrm{AI}}(t)$ and $D_{\mathrm{NW}}(t)$ are enumerated by brute force over all ancestral-independence families for $t\le3$ and, for every $t\le10$, every tuple of $D_{\mathrm{AI}}(t)$ is realized by an explicit family that the verifier checks against the definitions. The chain follows from the order-two and order-three brackets on the square, with the equality from Lemma~\ref{lem:rootsink}.
\end{proof}

The table shows two regularities that the theorems do not explain. The products $tq^H_t$ increase towards constants near $0.4$ on the square and $0.9$ on the triangle, an order of magnitude above the $1/16$ of Theorems~\ref{thm:squarelinear} and~\ref{thm:trianglelinear}; and the two hierarchies coincide along this direction at every odd order up to nine while separating at every even order, with a ratio $q^{\mathrm{NW}}_t/q^{\mathrm{AI}}_t-1$ decreasing like $0.4/t$. Neither pattern is proved beyond the certified orders; uncertified numerics at triangle order fifteen show a relative gap of $10^{-4}$ between the two hierarchies, which may be the end of the coincidence or conditioning error.
